\epanalipsi{5}

\tmima{Λεξικό του κεφαλαίου}

\renewcommand{\arraystretch}{1.5}
\begin{tabular*}{\linewidth}{@{\extracolsep{\fill}} >{\raggedright\arraybackslash}p{0.28\linewidth} >{\raggedright\arraybackslash}p{0.62\linewidth}}
\toprule
\textbf{Λέξη ή τύπος} & \textbf{Με μία πρόταση}\\
\midrule
$\alpha\%$ & Ο αριθμός $\dfrac{\alpha}{100}$: «$\alpha$ για κάθε $100$».\\
$\alpha\%$ του $\beta$ & $\dfrac{\alpha}{100}\cdot\beta$.\\
$\alpha\permil$ & Ο αριθμός $\dfrac{\alpha}{1\,000}$.\\
έκπτωση, αύξηση & Πολλαπλασιάζω με $1-\frac{\alpha}{100}$ ή $1+\frac{\alpha}{100}$.\\
ποσοστιαία μεταβολή & $\dfrac{\text{νέα}-\text{παλιά}}{\text{παλιά}}$, πάντα επί της αρχικής.\\
ΦΠΑ & Φόρος που προστίθεται στην αξία· μέσα σε τιμή με $24\%$ είναι τα $\frac{24}{124}$.\\
τόκος & $\text{κεφάλαιο}\cdot\text{επιτόκιο}\cdot\text{χρόνος}$· με κεφαλαιοποίηση ο τόκος προστίθεται στο κεφάλαιο.\\
\bottomrule
\end{tabular*}

\tmima{Ασκήσεις επανάληψης}

\ex Γράψε ως ποσοστά τα $\frac{3}{5}$, $\frac{7}{8}$ και τον $0{,}45$.

\ex Γράψε ως ανάγωγα κλάσματα τα $35\%$ και $12{,}5\%$.

\ex Βρες το $15\%$ των $240$ και το $60\%$ του $35$.

\ex Τι ποσοστό είναι τα $18$ από τα $72$;

\ex Μια τιμή των $140$ € έχει έκπτωση $30\%$. Πόσο κάνει;

\ex Μια τιμή των $250$ € αυξάνεται κατά $12\%$. Πόσο γίνεται;

\ex Πόσο \% αυξήθηκε μια τιμή από $80$ σε $100$;

\ex Μετά από έκπτωση $20\%$ ένα προϊόν κοστίζει $64$ €. Ποια ήταν η αρχική τιμή;

\ex Η τελική τιμή ενός προϊόντος με ΦΠΑ $24\%$ είναι $186$ €. Πόση είναι η αξία χωρίς ΦΠΑ;

\ex Κατάθεση $3\,000$ € με επιτόκιο $2\%$ τον χρόνο. Πόσος είναι ο τόκος ενός έτους και πόσος εννέα μηνών;

\ex Σε κυκλικό διάγραμμα $400$ ψήφων οι υποψήφιοι πήραν $45\%$, $30\%$, $15\%$ και ο τέταρτος το υπόλοιπο. Πόσο ποσοστό και πόσες ψήφους πήρε ο καθένας;

\ex Δύο διαδοχικές εκπτώσεις $10\%$ στα $200$ €. Πόσο κάνει και ποια είναι η συνολική έκπτωση;

\ex \textbf{Ερωτήσεις αυτοαξιολόγησης (ποσοστά).} Σωστό ή λάθος; α) $25\%=\frac{1}{4}$. \quad β) $0{,}07=7\%$. \quad γ) Το $150\%$ του $20$ είναι $30$. \quad δ) $10\permil=10\%$. \quad ε) Το $20\%$ του $50$ ισούται με το $50\%$ του $20$. \quad στ) $0{,}5\%=0{,}5$.

\ex \textbf{Ερωτήσεις αυτοαξιολόγησης (προβλήματα).} Σωστό ή λάθος; α) Έκπτωση $20\%$ και μετά $30\%$ κάνουν $50\%$. \quad β) Αν μια τιμή αυξηθεί κατά $25\%$ και μετά μειωθεί κατά $25\%$, επιστρέφει στην ίδια. \quad γ) Ο ΦΠΑ μέσα στην τιμή με $24\%$ είναι το $24\%$ της τελικής τιμής. \quad δ) Τα ποσοστά ενός πλήρους κυκλικού διαγράμματος κάνουν $100\%$. \quad ε) Ο τόκος ενός έτους είναι κεφάλαιο επί επιτόκιο.

\tmima{Ο χάρτης της επανάληψης}

Μετά τις λύσεις, δες σε ποιες ασκήσεις δυσκολεύτηκες. Ο πίνακας σου λέει ποια παράγραφο να ξαναδιαβάσεις.

\renewcommand{\arraystretch}{1.5}
\begin{tabular*}{\linewidth}{@{\extracolsep{\fill}} >{\raggedright\arraybackslash}p{0.38\linewidth} >{\raggedright\arraybackslash}p{0.52\linewidth}}
\toprule
\textbf{Αν δυσκολεύτηκες στις} & \textbf{ξαναδιάβασε τις παραγράφους}\\
\midrule
Ε5.1, Ε5.2, Ε5.3, Ε5.4, Ε5.13 & 5.1 (ποσοστά)\\
Ε5.5 έως Ε5.12, Ε5.14 & 5.2 (προβλήματα με ποσοστά)\\
\bottomrule
\end{tabular*}
